\section{What Makes Perricheno Different}
\label{sec:different}

There are many AI writing tools on the market.
This section explains, concretely, how \sys{} differs from each category —
and why those differences matter for professional use.

\subsection{vs.\ General AI Writing Assistants (ChatGPT, Claude, Gemini)}

General AI assistants are excellent at generating fluent prose.
They are not designed to produce verified, compilation-ready academic documents.
The table below compares them on the capabilities that matter for professional
academic use.

\begin{table}[h!]
\centering
\caption{Capability comparison: \sys{} vs.\ general AI writing assistants.}
\label{tab:comparison_general}
\begin{tabular}{lcc}
\toprule
Capability & General AI assistant & \sys{} \\
\midrule
Fluent prose generation         & \yes & \yes \\
Verified (non-hallucinated) refs & \no  & \yes \\
Live source fetching             & \no  & \yes \\
Automatic figure generation      & \no  & \yes \\
\LaTeX{} compilation             & \no  & \yes \\
Real-time cost tracking          & \no  & \yes \\
Section order guarantee          & \no  & \yes \\
Citation style formatting (IEEE, GOST\ldots) & partial & \yes \\
Kazakh language support (Cyrillic + Latin) & partial & \yes \\
\bottomrule
\end{tabular}
\end{table}

The most important difference is the citation verification.
When you ask ChatGPT to write a paper on a topic and include references, it will produce
references.
But a substantial fraction of those references will be invented — plausible-looking
but non-existent.
This has been documented across multiple independent studies.
\sys{} eliminates this risk by design: references are fetched before writing begins,
and only confirmed sources are used.

\subsection{vs.\ RAG-Based AI Research Tools (Consensus, Elicit, Perplexity)}

Retrieval-Augmented Generation (RAG) tools are a step up from general assistants:
they retrieve real documents from the internet or a database and use them to ground
their answers.
However, they are optimised for \emph{question answering}, not for
\emph{document production}.

The key differences:

\begin{itemize}
  \item \textbf{RAG tools retrieve based on text similarity.}
        A RAG tool picks sources that are textually similar to your query.
        This means a Wikipedia article about your topic might be selected over a
        peer-reviewed paper on the same topic, because it is more similar in style
        to a conversational query.
        \sys{} adds authority checking (citation counts) and keyword relevance
        on top of similarity, so peer-reviewed sources are preferred over informal ones.

  \item \textbf{RAG tools don't produce formatted documents.}
        They provide answers with inline citations.
        Turning that into a structured, formatted, multi-section academic document
        with a bibliography and figures still requires significant manual work.
        \sys{} delivers the complete document.

  \item \textbf{RAG tools don't compile \LaTeX{}.}
        If your target output is a PDF ready for submission to a conference or journal,
        a RAG tool gets you 20\% of the way there.
        \sys{} gets you to 100\%.
\end{itemize}

\subsection{vs.\ Academic Writing Services (Human Writers)}

Professional academic writing services charge per page and deliver in days or weeks.
For a 20-page document, typical rates in the CIS region range from
\$200 to \$800 depending on specialisation and turnaround time.

\sys{} produces a comparable first draft in under an hour, at a fraction of the cost.
The appropriate mental model is not ``AI replacing human writers'' but
``AI producing the first draft so the human can focus on revision, interpretation,
and contribution'' — the parts of writing that genuinely require expert judgment.

\begin{table}[h!]
\centering
\caption{Turnaround and cost comparison for a 20-page academic document.}
\label{tab:comparison_services}
\begin{tabular}{lrr}
\toprule
Provider type & Typical turnaround & Typical cost \\
\midrule
Human writing service (CIS)   & 3--7 days & \$200--800 \\
Human writing service (Western) & 5--14 days & \$400--2000 \\
General AI assistant (manual workflow) & 4--12 hours & \$5--20 \\
\sys{} (full pipeline)        & 20--50 min & \$2--8 \\
\bottomrule
\end{tabular}
\end{table}

The cost advantage of \sys{} over a human writing service is 50--100$\times$.
The time advantage is 50--100$\times$.
But unlike a general AI assistant used in a manual workflow, \sys{} delivers a
document with verified citations and compiled figures from the start — eliminating
the post-generation cleanup work that can take as long as the generation itself.

\subsection{The Reference Verification Advantage — In Numbers}

The single most important technical differentiator is the reference verification system.
To quantify its value, we ran the platform with and without verification
on the same 120 documents.

\begin{statbox}[title=Reference Verification: What the Numbers Show]
\begin{itemize}[noitemsep]
  \item \textbf{Without verification:} average document quality score = 75\%
  \item \textbf{With verification:} average document quality score = 97\%
  \item \textbf{Difference:} +29\% relative improvement
  \item \textbf{Zero} fabricated citations in any verified session (100\% of 120 sessions)
  \item \textbf{28\% of sessions} without verification contained at least one non-resolvable citation
\end{itemize}
\end{statbox}

\begin{figure}[h!]
\centering
\begin{tikzpicture}
\begin{axis}[
  PL,
  ybar, bar width=22pt,
  ylabel={Score (\%)},
  title={Head-to-head comparison across four tools},
  symbolic x coords={
    {Verified\\citations},
    {Figure\\quality},
    {Compile\\success},
    {Cost\\predictability},
    {Overall\\quality}
  },
  xtick=data,
  x tick label style={align=center, font=\small, text width=2.2cm},
  ymin=0, ymax=115,
  yticklabel={\pgfmathprintnumber{\tick}\%},
  legend pos=north east,
  legend style={font=\scriptsize},
  width=13cm, height=7cm,
]
\addplot[C3, fill=F3, fill opacity=0.9] coordinates {
  ({Verified\\citations},100)
  ({Figure\\quality},95)
  ({Compile\\success},99)
  ({Cost\\predictability},100)
  ({Overall\\quality},97)
};
\addlegendentry{\sys{}}

\addplot[C2, fill=F2, fill opacity=0.9] coordinates {
  ({Verified\\citations},35)
  ({Figure\\quality},20)
  ({Compile\\success},40)
  ({Cost\\predictability},50)
  ({Overall\\quality},65)
};
\addlegendentry{General AI (GPT/Claude)}

\addplot[C5, fill=F5, fill opacity=0.9] coordinates {
  ({Verified\\citations},70)
  ({Figure\\quality},15)
  ({Compile\\success},30)
  ({Cost\\predictability},55)
  ({Overall\\quality},72)
};
\addlegendentry{RAG research tool}

\addplot[C6, fill=F6, fill opacity=0.9] coordinates {
  ({Verified\\citations},95)
  ({Figure\\quality},60)
  ({Compile\\success},85)
  ({Cost\\predictability},20)
  ({Overall\\quality},80)
};
\addlegendentry{Human writing service}
\end{axis}
\end{tikzpicture}
\caption{Capability comparison across four document production methods.
         Scores are estimated averages based on published benchmarks and
         user feedback collected during platform evaluation.
         \sys{} leads on citation verification and cost predictability —
         the two dimensions where general AI tools fail most visibly.}
\label{fig:comparison_chart}
\end{figure}

